\documentclass[10pt,a4paper]{amsart}
\usepackage[margin=19mm]{geometry}
\usepackage{amsmath,amssymb,mathtools}
\usepackage[scr=rsfs,cal=euler]{mathalfa}
\usepackage{xfrac}
\usepackage[unicode,hidelinks,bookmarks=false]{hyperref}
\newcommand{\ie}{i.e.\ }
\newcommand{\cf}{cf.\ }
\newcommand{\resp}{respectively\ }
\newcommand{\Subsection}{\subsection}
\providecommand{\nobreakdash}{\nobreak\hbox{-}}
\setlength{\emergencystretch}{2em}
\multlinegap0pt
\usepackage{mathtools}
\usepackage{amssymb}
\usepackage{xcolor}
\usepackage{tikz}
\usepackage{enumerate}
\newtheorem{lem}{Lemma}
\newtheorem{prop}{Proposition}
\newcommand{\C}{\mathbb C}
\newcommand{\dd}{\,\mathrm d}
\title{Zero asymptotics for successive derivatives of hyperexponential functions with finite essential singularities}
\author{Christian H\"agg, Boris Shapiro}
\date{Source: arXiv:2604.09333v2 (CC BY 4.0)}
\begin{document}
\maketitle

\noindent\emph{Source and licence.} Zero asymptotics for successive derivatives of hyperexponential functions with finite essential singularities. Version arXiv:2604.09333v2 (CC BY 4.0), distributed by arXiv under the Creative Commons Attribution 4.0 International licence, http://creativecommons.org/licenses/by/4.0/, as recorded in the arXiv licence metadata for this exact version and shown on the abstract page https://arxiv.org/abs/2604.09333v2. Full licence notice: \texttt{licenses/arxiv-2604.09333-CC-BY-4.0.txt}.\par\smallskip

\noindent\textit{Editorial scope.} Counted unit: the article's lem lem:wright-multi (source0.tex lines 711--757). The article's complete original proof of it is delivered with it (source0.tex lines 759--858, one \texttt{proof} environment of 100 lines). No mathematics is written, repaired or re-derived: the statement and the proof are the article's own lines, in the article's own notation, and no retained line is substituted or re-flowed.  Retained uncounted as the source-local context the proof needs: lines 509--513; lines 521--544; lines 582--594; lines 623--640; lines 860--895.  Nothing was omitted from any retained span, so the package carries no omission record.  No retained line contains a citation, so the wrapper carries no bibliography and no bibliography entry.  No retained line contains a figure environment, an image-inclusion command or drawing code, so no image is delivered and the package declares no asset dependency.  Editorial formatting only, in addition: an AMS \texttt{amsart} preamble that restates the article's own declarations and macros used by the retained lines, namely the environments \texttt{lem}, \texttt{prop} on separate counters, together with the standard packages those lines need.  The retained environments print in this document as Lemma 1, Proposition 1 in order of appearance, whereas the article prints the same items with the numbering of its own full document.  The complete native source of the article as distributed in the arXiv e-print for this exact version is bundled unchanged as \texttt{sources/198202/source0.tex}.
\begin{lem}[landing and continuation of level-gradient separatrices]\label{lem:gradient-landing}
Let $X$ be a compact bordered Riemann surface with piecewise $C^1$ boundary, let $\Phi=h+ig$ be holomorphic on a neighborhood of $X$, and assume that all critical points of $\Phi$ in $X^\circ$ are simple and have distinct $g$-values. Let $A\cup B=\partial X$ be a decomposition into finitely many closed boundary arcs, meeting only at endpoints, and let $I\subset\mathbb R$ be a compact interval which contains the $g$-values of the critical points and is disjoint from the $g$-values of the endpoints $A\cap B$. Suppose that, on $\partial X\cap g^{-1}(I)$, the field $-\nabla h$ has normal component bounded away from zero, points outward on the relative interiors of the $B$-arcs, and points inward on the relative interiors of the $A$-arcs. Then every downward separatrix from a critical point hits the relative interior of $B$ in finite time and transversely, and no downward saddle connection occurs. The corresponding upward statement holds with $A$ and $B$ interchanged.

If the data depend $C^1$-continuously on a parameter in a compact set, the boundary decomposition is locally trivial, and the preceding hypotheses hold with uniform margins, then the landed separatrix arcs continue by ambient isotopy under parameter variation. In particular their relative homology classes, and the corresponding intersection numbers with a fixed relative cycle transported by the same boundary trivialization, are locally constant.
\end{lem}

\begin{lem}[uniform annular exit and flow data]\label{lem:wright-exit-arcs}
Let $V$ be a connected component of $D\setminus\mathfrak S$, and let $U_0$ be connected with $\overline{U_0}\Subset V$. Put
$$\tau_\nu(z):=\omega_\nu\eta(z),\qquad \psi_z(\tau):=\lambda_m(z)\tau^{-m}+\tau.$$
Choose $\Lambda>0$ such that $|\tau_\nu(z)|<\Lambda/4$ on $\overline{U_0}$ for every $\nu$. Then there are constants $C_0>0$, $0<\delta<\Lambda<\Lambda_1$, $q_0>0$, and $b>0$, with $q_0\Lambda_1<1$, for which the following assertions hold. Set
$$X:=\{\delta\le |\tau|\le\Lambda_1\},\qquad \mathcal P:=\overline{U_0}\times[0,q_0].$$
For $0<q\le q_0$ define
$$\widehat\Phi_{z,q}(\tau):=\sum_{s=1}^m q^{m-s}\lambda_s(z)\tau^{-s}-(q^{-1}+q^m)\log(1-q\tau),$$
using the Taylor branch of the logarithm on $X$, and set $\widehat\Phi_{z,0}:=\psi_z$. Write
$$h_{z,q}:=\Re\widehat\Phi_{z,q},\qquad g_{z,q}:=\Im\widehat\Phi_{z,q},$$
and let $\mathbf n$ be the outward unit normal of the annulus $X$. For every $(z,q)\in\mathcal P$ the set
$$
B_{z,q}:=\overline{\{\tau\in\partial X:\ \partial_{\mathbf n}h_{z,q}(\tau)<0,\ |g_{z,q}(\tau)|\le C_0+1\}}
$$
has exactly $m+1$ connected components: $m$ compact intervals on $|\tau|=\delta$ and one compact interval on $|\tau|=\Lambda_1$. Put
$$A_{z,q}:=\overline{\partial X\setminus B_{z,q}}.$$
Then $A_{z,q}\cup B_{z,q}=\partial X$, $A_{z,q}\cap B_{z,q}$ is the finite set of endpoints of these intervals, and $A_{z,q}$ also has $m+1$ interval components. The endpoints are transverse intersections with the levels $g_{z,q}=\pm(C_0+1)$, the pairs $(A_{z,q},B_{z,q})$ vary locally trivially in $(z,q)$ as boundary subcomplexes, and no boundary tangency of $\pm\nabla h_{z,q}$ occurs on $\partial X\cap\{|g_{z,q}|\le C_0+1\}$. On $\partial X\cap\{|g_{z,q}|\le C_0\}$, the downward field $-\nabla h_{z,q}$ points strictly outward exactly on the relative interiors of $B_{z,q}$ and strictly inward on the relative interiors of $A_{z,q}$; the upward field has the opposite signs.

For every $(z,q)\in\mathcal P$, the phase $\widehat\Phi_{z,q}$ has exactly $m+1$ simple critical points in $X$, labelled by
$$\tau_\nu(z,q)=\tau_\nu(z)+O_{U_0}(q),\qquad 0\le\nu\le m,$$
with $\tau_\nu(z,0)=\tau_\nu(z)$. Their critical imaginary parts are separated by a positive constant independent of $(z,q)$, and $|g_{z,q}(\tau_\nu(z,q))|\le C_0$. Every maximal downward separatrix from a critical point hits the relative interior of $B_{z,q}$ in finite time and transversely, and no downward trajectory joins two critical points. Every maximal upward separatrix from a critical point hits the relative interior of $A_{z,q}$ in finite time and transversely, and no upward trajectory joins two critical points. Moreover,
$$\sup_{\tau\in B_{z,q}}\Re\widehat\Phi_{z,q}(\tau)
\le \min_{0\le\nu\le m}\Re\widehat\Phi_{z,q}(\tau_\nu(z,q))-b$$
uniformly on $\mathcal P$.
\end{lem}

\begin{lem}[relative thimbles and contour coefficients]\label{lem:wright-thimble-basis}
Assume the notation and conclusions of Lemma~\ref{lem:wright-exit-arcs}. For each $(z,q)\in\mathcal P$, let $\Delta_\nu(z,q)$ be the downward thimble through $\tau_\nu(z,q)$, viewed as a relative one-cycle in $H_1(X,B_{z,q};\mathbb Z)$, and let $\nabla_\nu(z,q)$ be the upward thimble through the same saddle, viewed in $H_1(X,A_{z,q};\mathbb Z)$. The relative groups are those of the boundary triad
$$A_{z,q}\cup B_{z,q}=\partial X,\qquad A_{z,q}\cap B_{z,q}\text{ finite}.$$
Choose an orientation of each downward thimble and orient the corresponding upward thimble so that the local intersection at the common saddle is $+1$; for the thimble-thimble pairing, perturb boundary endpoints only inside boundary collars when needed. Then
$$\Delta_\nu(z,q)\cdot\nabla_\mu(z,q)=\delta_{\nu\mu}.$$
The classes $[\Delta_\nu(z,q)]$ form an integral basis of $H_1(X,B_{z,q};\mathbb Z)$, and $[\nabla_\nu(z,q)]$ form the dual integral basis of $H_1(X,A_{z,q};\mathbb Z)$. Hence, for the positively oriented circle $\Gamma_\Lambda:=\{|\tau|=\Lambda\}$,
$$[\Gamma_\Lambda]=\sum_{\nu=0}^m\sigma_\nu[\Delta_\nu(z,q)]
\quad\text{in }H_1(X,B_{z,q};\mathbb Z),$$
where
$$\sigma_\nu=[\Gamma_\Lambda]\cdot[\nabla_\nu(z,q)]\in\mathbb Z$$
is the homological intersection pairing. If this number is evaluated by crossing counts, use any sufficiently small transverse perturbation of the core representative inside $X$.
After these thimble orientations are transported through the locally trivial family of boundary triads, the integers $\sigma_\nu$ are independent of $(z,q)\in\mathcal P$, and they are not all zero.
\end{lem}

\begin{lem}[filtered deformation for the scaled Wright phase]\label{lem:filtered-wright}
Let $V$ be a connected component of $D\setminus\mathfrak S$, let $U_0$ be connected with $\overline{U_0}\Subset V$, and let $K_0\Subset U_0$. Put
$$\tau_\nu(z):=\omega_\nu\eta(z),\qquad \psi_z(\tau):=\lambda_m(z)\tau^{-m}+\tau.$$
Choose $\Lambda>0$ such that $|\tau_\nu(z)|<\Lambda/4$ on $\overline{U_0}$ for every $\nu$. Then there are $0<\delta<\Lambda<\Lambda_1$ and $q_0>0$, with $q_0\Lambda_1<1$, such that Lemmas~\ref{lem:wright-exit-arcs} and~\ref{lem:wright-thimble-basis} apply to $\widehat\Phi_{z,q}$ on $X=\{\delta\le |\tau|\le\Lambda_1\}$. Let
$$I:=\{\nu:\sigma_\nu\ne0\},\qquad
\mathcal H_{z,q}:=\max_{\nu\in I}\Re\widehat\Phi_{z,q}(\tau_\nu(z,q)).$$
There are pairwise disjoint neighborhoods $\mathcal N_\mu$ of the critical graphs $\{(z,q,\tau_\mu(z,q)):(z,q)\in K_0\times[0,q_0]\}$, each carrying a parameter-dependent holomorphic Morse coordinate for $\widehat\Phi_{z,q}$, and a number $c>0$ such that
$$C_{z,q}:=\sum_{\nu\in I}\sigma_\nu\Delta_\nu(z,q)$$
represents $[\Gamma_\Lambda]$ in $H_1(X,B_{z,q};\mathbb Z)$, contains no zero-coefficient thimble at chain level, has uniformly bounded piecewise $C^1$ length and multiplicity, and satisfies
$$\operatorname{supp}C_{z,q}\setminus\bigcup_{\nu\in I}\mathcal N_\nu(z,q)
\subset\{\tau\in X:\Re\widehat\Phi_{z,q}(\tau)\le\mathcal H_{z,q}-c\}.$$
Moreover one may choose singular two-chains $\mathcal S_{z,q}$ in $X$ and one-chains $E_{z,q}$ in $B_{z,q}$ such that
$$\Gamma_\Lambda-C_{z,q}=\partial\mathcal S_{z,q}+E_{z,q},\qquad
\partial E_{z,q}=-\partial C_{z,q},$$
$$\operatorname{supp}E_{z,q}\subset\{\tau\in B_{z,q}:\Re\widehat\Phi_{z,q}(\tau)\le\mathcal H_{z,q}-c\},$$
and $E_{z,q}$ has uniformly bounded piecewise $C^1$ length and multiplicity. Consequently, for every holomorphic one-form $\Omega$ on a neighborhood of $X$,
$$\int_{\Gamma_\Lambda}\Omega=\int_{C_{z,q}}\Omega+\int_{E_{z,q}}\Omega.$$
\end{lem}

\begin{prop}[active Wright coefficients]\label{prop:active-wright-coefficients}
Let $V$ be a leading Stokes chamber in the notation of this subsection, and fix $z_\ast\in V$. Put
\[
\psi_\ast(\tau):=\lambda_m(z_\ast)\tau^{-m}+\tau,
\qquad
\tau_{\nu,\ast}:=\omega_\nu\eta(z_\ast),\quad 0\le\nu\le m.
\]
Set $v_\nu:=\omega_\nu\eta(z_\ast)$ and $P_\ast:=\operatorname{conv}\{v_0,\ldots,v_m\}$. The limiting critical values of $\psi_\ast$ at the saddles $\tau_{\nu,\ast}$ are $\frac{m+1}{m}v_\nu$, so, up to this common factor, $P_\ast$ is the regular $(m+1)$-gon of limiting critical values for the Wright phase $\psi_\ast$. Since $z_\ast\notin\mathfrak S$, no two $v_\nu$ have the same imaginary part. Let $\mathcal R_V$ be the set of labels for the vertices on the boundary chain of $P_\ast$ visible from $+\infty$ in the real direction. Equivalently, with $\operatorname{Arg}\in(-\pi,\pi]$,
\begin{equation}\label{eq:active-right-boundary}
\mathcal R_V=
\left\{0\le\nu\le m:
\left|\operatorname{Arg} v_\nu\right|<
\frac{\pi}{2}+\frac{\pi}{m+1}\right\}.
\end{equation}
Equality in \eqref{eq:active-right-boundary} is equivalent to equality of the imaginary parts of two adjacent limiting critical values. These adjacent equalities are among the components of \(\mathfrak S\), so equality cannot occur for \(z_\ast\in V\). Hence \(\mathcal R_V\) is independent of the base point in \(V\). Non-adjacent components of \(\mathfrak S\) may further subdivide the parameter domain, but they do not change the visible-chain criterion. The coefficient-contour coordinates in Lemmas~\ref{lem:wright-thimble-basis} and~\ref{lem:wright-multi} satisfy
\[
|\sigma_{\nu,V}|=
\begin{cases}
1, & \nu\in\mathcal R_V,\\
0, & \nu\notin\mathcal R_V.
\end{cases}
\]
Thus
\[
I_V=\{\nu:\sigma_{\nu,V}\ne0\}=\mathcal R_V.
\]
The signed value for $\nu\in\mathcal R_V$ is the oriented crossing sign of the corresponding upward thimble with the positively oriented core circle, with the Morse orientation convention of Lemma~\ref{lem:wright-multi}. These signs are independent of the admissible annulus, core representative, and base point in $V$, after continuation of the saddle labels and Morse orientations.

For an essential site $a_i$, the same rule is applied with
\[
v_\nu=\omega_\nu\eta_D(z_\ast),
\qquad
\eta_D(z_\ast)^{m_i+1}=m_i\lambda_{i,m_i}(z_\ast-a_i)^{-m_i}.
\]
Changing the branch of $\eta_D$ only relabels the output.
\end{prop}

\begin{lem}[sectorial parameter-uniform multi-saddle asymptotics]\label{lem:wright-multi}
Let $V$ be a connected component of $D\setminus\mathfrak S$, and let $K\Subset V$ be compact. Then there exists a connected open set $U$ with $K\Subset U\Subset V$ such that, for each $0\le\nu\le m$ and all sufficiently large $n$, there is a holomorphic function $t_{n,\nu}:U\to\C$ satisfying
$$\Phi_{z,n}'(t_{n,\nu}(z))=0,\qquad
t_{n,\nu}(z)=\omega_\nu\eta(z)\,n^{-1/(m+1)}+O_U\!\left(n^{-2/(m+1)}\right).$$
Set
\begin{equation}\label{eq:wright-multi-data}
\theta:=-\frac{2\beta+m+2}{2(m+1)},\qquad
A_\nu(z):=\frac{(\omega_\nu\eta(z))^{\beta}\,(\omega_\nu^{1/2}\eta^{1/2}(z))}{\sqrt{2\pi(m+1)}}\,R(z,1),
\qquad
\Xi_\nu(z;n):=\Phi_{z,n}(t_{n,\nu}(z)).
\end{equation}
Then each $A_\nu$ is holomorphic and nonvanishing on $U$, and each $\Xi_\nu(z;n)$ admits, by iterating the saddle equation, an asymptotic expansion to arbitrary finite order in descending powers of $n^{1/(m+1)}$, uniformly on $K$, with
\begin{equation}\label{eq:wright-multi-phase}
\Xi_\nu(z;n)=\frac{m+1}{m}\,\omega_\nu\eta(z)\,n^{m/(m+1)}
+O_K\!\left(n^{(m-1)/(m+1)}\right)
\end{equation}
uniformly for $z\in K$.

Write $\tau_\nu(z):=\omega_\nu\eta(z)$ and use the square roots in \eqref{eq:wright-multi-data} to set $\tau_\nu(z)^{1/2}=\omega_\nu^{1/2}\eta^{1/2}(z)$. Orient each local downward thimble at $t_{n,\nu}(z)$ so that, in the Morse coordinate $u$ normalized by
$$\Phi_{z,n}(t)=\Xi_\nu(z;n)-u^2/2,
\qquad
\left.\frac{\partial t}{\partial u}\right|_{u=0}=\frac{i}{\sqrt{\Phi_{z,n}''(t_{n,\nu}(z))}},$$
the real $u$-axis has its standard orientation; the square root of $\Phi_{z,n}''$ is the continuation whose leading term is
$$\sqrt{\frac{m+1}{\tau_\nu(z)}}\,n^{(m+2)/(2(m+1))}.$$
With these local thimble orientations, Lemma~\ref{lem:filtered-wright}, applied on $U\Subset V$ to the scaled coefficient circle, gives integers $\sigma_{\nu,U}$, not all zero. Because the branch of $\eta$ has been fixed on $D$, the limiting saddle labels $\tau_\nu(z)=\omega_\nu\eta(z)$ are single-valued on all of $V$; for large $n$ the finite-$q$ saddles are labelled by the perturbation of these branches. If $U_1,U_2\Subset V$ are two admissible choices, choose a compact connected set $L\Subset V$ containing $U_1\cup U_2$ and cover $L$ by finitely many subdomains to which Lemmas~\ref{lem:wright-exit-arcs}--\ref{lem:filtered-wright} apply with overlapping closures. On each overlap the locally trivial boundary triads identify the labelled thimbles, and Lemma~\ref{lem:gradient-landing} excludes corner hits, boundary tangencies, critical collisions and saddle connections during the continuation. The intersection numbers with the transported coefficient-core class therefore agree from one member of the finite chain to the next. Thus the integers are independent of the auxiliary compact chart and are denoted by $\sigma_{\nu,V}$; this notation records homological continuation on compact subsets of $V$ and does not assert a uniform finite-$q$ isotopy on all of $V$ at once. Put
$$I_V:=\{\nu:\sigma_{\nu,V}\ne0\}.$$
Then $I_V$ is nonempty. Proposition~\ref{prop:active-wright-coefficients} below computes the vanishing pattern of these integers explicitly; the present lemma uses only their homological definition and nontriviality. If
$$M_V(z;n):=\max_{\nu\in I_V}\Re\Xi_\nu(z;n),$$
then there are functions $\varepsilon_{n,\nu}:K\to\C$, $\nu\in I_V$, satisfying
$$\max_{\nu\in I_V}\sup_{z\in K}|\varepsilon_{n,\nu}(z)|=O_K\!\left(n^{-1/(m+1)}\right),$$
and, uniformly for $z\in K$,
\begin{equation}\label{eq:wright-multi}
[\zeta^n]F_z(\zeta)
=
n^\theta\sum_{\nu\in I_V}\sigma_{\nu,V}A_\nu(z)\exp\!\bigl(\Xi_\nu(z;n)\bigr)\bigl(1+\varepsilon_{n,\nu}(z)\bigr)
+\mathcal R_{n,V}(z),
\end{equation}
with
\begin{equation}\label{eq:wright-multi-rem}
\mathcal R_{n,V}(z)=
O_K\!\left(
n^\theta
\exp\!\bigl(M_V(z;n)-c_K n^{m/(m+1)}\bigr)
\right)
\end{equation}
for some $c_K>0$. The estimate is uniform on compact sets crossing leading anti-Stokes arcs. Changing one square-root branch in \eqref{eq:wright-multi-data} changes the displayed $A_\nu$ by a sign; changing the corresponding thimble orientation changes $\sigma_{\nu,V}$ by the same sign, so the products $\sigma_{\nu,V}A_\nu$ are intrinsic.
\end{lem}

\begin{proof}
Fix $K\Subset V$ and choose a connected open set $U$ with $K\Subset U\Subset V$. Put $\tau_\nu(z):=\omega_\nu\eta(z)$. The scaled saddle equation is
\begin{align*}
0
&=n^{-1}\Phi_{z,n}'\!\bigl(n^{-1/(m+1)}\tau\bigr)\\
&=1-\tau_\nu(z)^{m+1}\tau^{-m-1}+O_U\!\left(n^{-1/(m+1)}\right),
\end{align*}
uniformly for $\tau$ on compact subsets of $\C^\times$. Since
$$\partial_\tau\!\left(1-\tau_\nu(z)^{m+1}\tau^{-m-1}\right)\Big|_{\tau=\tau_\nu(z)}=\frac{m+1}{\tau_\nu(z)}\ne0,$$
the holomorphic implicit function theorem gives holomorphic saddles $t_{n,\nu}=n^{-1/(m+1)}\tau_{n,\nu}$ on $U$, with
$$\tau_{n,\nu}(z)=\tau_\nu(z)+O_U\!\left(n^{-1/(m+1)}\right),
\qquad
 t_{n,\nu}(z)=\tau_\nu(z)n^{-1/(m+1)}+O_U\!\left(n^{-2/(m+1)}\right).$$
Repeated substitution in the implicit equation gives the full expansion of $\tau_{n,\nu}$ and then of $\Xi_\nu$. Its leading term is
$$\lambda_m(z)\tau_\nu(z)^{-m}n^{m/(m+1)}+\tau_\nu(z)n^{m/(m+1)}
=\frac{m+1}{m}\,\tau_\nu(z)n^{m/(m+1)},$$
which is \eqref{eq:wright-multi-phase}.

By the annulus hypothesis applied to $\overline U$, choose $\varrho>1$ such that $(z,\zeta)\mapsto F_z(\zeta)$ is holomorphic on a neighborhood of $\overline U\times(\{|\zeta|\le\varrho\}\setminus\{1\})$. Choose $\Lambda>0$ so that every $\tau_\nu(z)$ lies in $|\tau|<\Lambda/4$ for $z\in\overline U$, and put $\varepsilon_n:=\Lambda n^{-1/(m+1)}$. For $0<r<1$ with $|1-\zeta|>\varepsilon_n$ on $|\zeta|=r$, Cauchy's formula gives
$$[\zeta^n]F_z=(2\pi i)^{-1}\int_{|\zeta|=r}F_z(\zeta)\zeta^{-n-1}\,\dd\zeta.$$
Cauchy's theorem in
$$\{r\le |\zeta|\le\varrho\}\setminus\{|1-\zeta|<\varepsilon_n\}$$
gives, with the boundary orientations induced by this punctured annulus,
\begin{equation}\label{eq:wright-contour-split}
[\zeta^n]F_z(\zeta)=\frac{1}{2\pi i}\int_{|\zeta|=\varrho}F_z(\zeta)\zeta^{-n-1}\,\dd\zeta
+\frac{1}{2\pi i}\int_{\Gamma_{\varepsilon_n}} t^\beta R(z,1-t)e^{\Phi_{z,n}(t)}\,\dd t,
\end{equation}
where $t=1-\zeta$ and $\Gamma_{\varepsilon_n}$ is the positively oriented circle $|t|=\varepsilon_n$. The sign is positive because the small boundary is clockwise in the $\zeta$-plane and $\dd\zeta=-\dd t$. The outer integral is $O_U(\varrho^{-n})$.

Set $q_n:=n^{-1/(m+1)}$ and
$$\widehat\Phi_{z,n}(\tau):=n^{-m/(m+1)}\Phi_{z,n}\!\bigl(q_n\tau\bigr).$$
Apply Lemma~\ref{lem:filtered-wright} to $U_0=U$, $K_0=K$, and this $\Lambda$. For all large $n$, $q_n\le q_0$ and
$$\widehat\Phi_{z,q_n}(\tau)=\widehat\Phi_{z,n}(\tau),$$
so the labelled critical points in the scaled annulus are $\tau_{n,\nu}(z)=q_n^{-1}t_{n,\nu}(z)$. The scaled image of $\Gamma_{\varepsilon_n}$ is the positively oriented circle $\Gamma_\Lambda:=\{|\tau|=\Lambda\}$. Lemma~\ref{lem:wright-thimble-basis} gives
$$[\Gamma_\Lambda]=\sum_{\nu=0}^m\sigma_{\nu,V}[\Delta_\nu(z,q_n)]
\quad\text{in }H_1(X,B_{z,q_n};\mathbb Z),$$
where the integers are locally constant on the compact parameter family $\overline U\times[0,q_0]$. More explicitly, for fixed $U\Subset V$ the path $q\in[0,q_0]$ stays inside the family of Lemma~\ref{lem:wright-exit-arcs}: the boundary triad is locally trivial, critical points neither collide nor acquire equal imaginary values, and Lemma~\ref{lem:gradient-landing} excludes creation of saddle connections along the path. Therefore the relative homology local system over $\overline U\times[0,q_0]$ is trivialized by the transported thimbles, and the intersection number $[\Gamma_\Lambda]\cdot[\nabla_\nu(z,q)]$ is independent of $q$. Thus the coefficients used at $q_n=n^{-1/(m+1)}$ for all sufficiently large $n$ are the same intersection numbers as for the limiting Wright phase $\lambda_m(z)\tau^{-m}+\tau$ on $U$. The notation $\sigma_{\nu,V}$ records the continuation of these limiting coefficients in the chamber; no uniform assertion is used for $q$-dependent points approaching $\partial V$. Hence $I_V:=\{\nu:\sigma_{\nu,V}\ne0\}$ is nonempty.

Let
$$H_{z,n}:=\max_{\nu\in I_V}\Re\widehat\Phi_{z,n}(\tau_{n,\nu}(z))=n^{-m/(m+1)}M_V(z;n).$$
Lemma~\ref{lem:filtered-wright} gives a filtered chain $C_{z,q_n}$ and a boundary chain $E_{z,n}$ such that
$$\int_{\Gamma_\Lambda}\Omega_{z,n}=\int_{C_{z,q_n}}\Omega_{z,n}+\int_{E_{z,n}}\Omega_{z,n},$$
where
$$\Omega_{z,n}:=q_n^{\beta+1}\tau^\beta R(z,1-q_n\tau)
\exp\!\bigl(n^{m/(m+1)}\widehat\Phi_{z,n}(\tau)\bigr)\,\dd\tau.$$
This one-form is holomorphic on a neighborhood of $X$ for all large $n$: $\beta\in\mathbb Z$, $0\notin X$, $q_n\Lambda_1<1$ keeps the logarithmic singularity outside the annulus and fixes the Taylor branch, and $R(z,1-q_n\tau)$ is holomorphic there. Since $q_n\Lambda_1\to0$, the factor $R(z,1-q_n\tau)$ and all derivatives needed below are bounded uniformly for $z\in K$ and $\tau\in X$. Moreover,
$$\operatorname{supp}C_{z,q_n}\setminus\bigcup_{\nu\in I_V}\mathcal N_\nu(z,q_n)
\subset\{\Re\widehat\Phi_{z,n}\le H_{z,n}-c\},
\qquad
\operatorname{supp}E_{z,n}\subset\{\Re\widehat\Phi_{z,n}\le H_{z,n}-c\}$$
uniformly for $z\in K$ and all large $n$. Thus inactive thimbles are absent from the chain representative, and the boundary chain is included in the integral identity.

On the nonlocal pieces of $C_{z,q_n}$ and on all of $E_{z,n}$,
$$\Re\Phi_{z,n}\le M_V(z;n)-c n^{m/(m+1)}.$$
The length and multiplicity of these chains are uniformly bounded. The amplitude $q_n^{\beta+1}\tau^\beta R(z,1-q_n\tau)$ is bounded by a fixed power of $n$ on $X$. Decreasing $c$ absorbs this polynomial factor and gives the remainder estimate \eqref{eq:wright-multi-rem} for all nonlocal and boundary pieces. Since $n^{-m/(m+1)}\Xi_\nu(z;n)$ is uniformly bounded on $K$, $M_V(z;n)\ge -C_Kn^{m/(m+1)}$; the outer integral $O_U(\varrho^{-n})$ in \eqref{eq:wright-contour-split} is then also absorbed by \eqref{eq:wright-multi-rem}, after decreasing $c$ once more.

It remains to evaluate the active saddle neighborhoods. Fix $\nu\in I_V$, put $\alpha=m/(m+1)$, and write $\tau_s=\tau_{n,\nu}(z)$ and $t_s=q_n\tau_s$. The parametric holomorphic Morse lemma, applied on the compact critical graph over $K$, gives one radius $r>0$ and coordinates $v$ with uniformly bounded inverse maps and derivatives such that
$$\widehat\Phi_{z,n}(\tau)=\widehat\Phi_{z,n}(\tau_s)-v^2/2,
\qquad
\left.\frac{\partial\tau}{\partial v}\right|_{v=0}=\frac{i}{\sqrt{\widehat\Phi_{z,n}''(\tau_s)}},$$
and the oriented local downward thimble is the real $v$-axis. The square root is chosen so that
$$\widehat\Phi_{z,n}''(\tau_s)=\frac{m+1}{\tau_\nu(z)}\left(1+O_U\!\left(n^{-1/(m+1)}\right)\right),$$
which is equivalent to the convention
$$\sqrt{\Phi_{z,n}''(t_s)}=
\sqrt{\frac{m+1}{\tau_\nu(z)}}\,n^{(m+2)/(2(m+1))}
\left(1+O_U\!\left(n^{-1/(m+1)}\right)\right).$$
Define
$$a_{z,n}(v):=\tau(v)^\beta R(z,1-q_n\tau(v))\frac{\partial\tau}{\partial v}(v).$$
The functions $a_{z,n}$ have uniformly bounded first two derivatives for $|v|\le r$. The local contribution of the unit oriented thimble is
\begin{align*}
\frac{1}{2\pi i}\int q_n^{\beta+1}\tau^\beta R(z,1-q_n\tau)
 e^{n^\alpha\widehat\Phi_{z,n}(\tau)}\,\dd\tau
&=\frac{e^{\Xi_\nu(z;n)}q_n^{\beta+1}}{2\pi i}
\int_{-r}^r a_{z,n}(v)e^{-n^\alpha v^2/2}\,\dd v\\
&\quad +O_K\!\left(n^A e^{\Re\Xi_\nu(z;n)-c n^\alpha}\right)
\end{align*}
for a fixed $A$. The error includes the part of the thimble outside $|v|<r$ and the Gaussian tail, both of which lie a fixed amount below the saddle height in the scaled phase. Taylor expansion at $v=0$ gives
$$\int_{-r}^r a_{z,n}(v)e^{-n^\alpha v^2/2}\,\dd v
=a_{z,n}(0)\sqrt{2\pi}\,n^{-\alpha/2}
\left(1+O_K\!\left(n^{-1/(m+1)}\right)\right),$$
because the odd term integrates to zero and the next even term is $O(n^{-\alpha})$, which is $O(n^{-1/(m+1)})$ for $m\ge1$.

Since $a_{z,n}(0)=\tau_s^\beta R(z,1-t_s)i/\sqrt{\widehat\Phi_{z,n}''(\tau_s)}$ and
$$\Phi_{z,n}''(t_s)=n^\alpha q_n^{-2}\widehat\Phi_{z,n}''(\tau_s),$$
the preceding display becomes
$$\exp\!\bigl(\Xi_\nu(z;n)\bigr)
\frac{t_s^\beta R(z,1-t_s)}{\sqrt{2\pi\Phi_{z,n}''(t_s)}}
\left(1+O_K\!\left(n^{-1/(m+1)}\right)\right).$$
Finally,
$$t_s^\beta R(z,1-t_s)=q_n^\beta\tau_\nu(z)^\beta R(z,1)
\left(1+O_K\!\left(n^{-1/(m+1)}\right)\right),$$
and the square-root convention gives
$$\frac{1}{\sqrt{\Phi_{z,n}''(t_s)}}=
\frac{\tau_\nu(z)^{1/2}}{\sqrt{m+1}}\,n^{-(m+2)/(2(m+1))}
\left(1+O_K\!\left(n^{-1/(m+1)}\right)\right).$$
Thus
$$\frac{t_s^\beta R(z,1-t_s)}{\sqrt{2\pi\Phi_{z,n}''(t_s)}}
=n^\theta A_\nu(z)\left(1+O_K\!\left(n^{-1/(m+1)}\right)\right),$$
with the $A_\nu$ and $\theta$ of \eqref{eq:wright-multi-data}. Multiplying these unit-thimble contributions by the fixed integers $\sigma_{\nu,V}$, summing over $\nu\in I_V$, and adding the estimated remainder proves \eqref{eq:wright-multi}.
\end{proof}

\end{document}
